% ------------------------------------------------------------------
% Gemeinsame Praeambel fuer main.tex (Aufgaben) und loesungen.tex.
% Die \documentclass-Zeile steht in main.tex / loesungen.tex,
% weil sie sich nur in solution=false/true unterscheidet.
% ------------------------------------------------------------------
\usepackage[english, main=ngerman]{babel}
\usepackage[utf8]{inputenc}
\usepackage[autostyle]{csquotes}
\usepackage{listings}
\definecolor{mauve}{rgb}{0.58,0,0.82}   % fuer stringstyle unten; war nicht definiert

\lstloadlanguages{C}
\lstset{commentstyle=\itshape, stringstyle=\sffamily}
\lstset{inputencoding=utf8, language=C, frame=none, basicstyle=\ttfamily, keepspaces=true, extendedchars=true, tabsize=4, columns=flexible, inputencoding=utf8,
    extendedchars=true,
    literate={ä}{{\"a}}1 {ö}{{\"o}}1 {ü}{{\"u}}1 {Ä}{{\"A}}1 {Ö}{{\"O}}1 {Ü}{{\"U}}1 }
\lstset{escapechar=\%}

\lstdefinelanguage[RISC-V]{Assembler}
{
  alsoletter={.}, % allow dots in keywords
  alsodigit={0x}, % hex numbers are numbers too!
  morekeywords=[1]{ % instructions
    lb, lh, lw, lbu, lhu,
    sb, sh, sw,
    sll, slli, srl, srli, sra, srai,
    add, addi, sub, lui, auipc, li, mv,
    mul, mulh, mulhsu, mulhu, div, divu, rem, remu,
    xor, xori, or, ori, and, andi,
    slt, slti, sltu, sltiu,
    beq, bne, blt, bge, bltu, bgeu,
    j, jr, jal, jalr, ret,
    scall, break, nop
  },
  morekeywords=[2]{ % sections of our code and other directives
    .align, .ascii, .asciiz, .byte, .data, .double, .extern,
    .float, .globl, .half, .kdata, .ktext, .set, .space, .text, .word
  },
  morekeywords=[3]{ % registers
    zero, ra, sp, gp, tp, s0, fp,
    t0, t1, t2, t3, t4, t5, t6,
    s1, s2, s3, s4, s5, s6, s7, s8, s9, s10, s11,
    a0, a1, a2, a3, a4, a5, a6, a7,
    ft0, ft1, ft2, ft3, ft4, ft5, ft6, ft7,
    fs0, fs1, fs2, fs3, fs4, fs5, fs6, fs7, fs8, fs9, fs10, fs11,
    fa0, fa1, fa2, fa3, fa4, fa5, fa6, fa7
  },
  morecomment=[l]{;},   % mark ; as line comment start
  morecomment=[l]{\#},  % as well as # (even though it is unconventional)
  morestring=[b]",      % mark " as string start/end
  morestring=[b]'       % also mark ' as string start/end
}

\lstset{
  % listings sonderzeichen (for german weirdness)
  literate={ö}{{\"o}}1
           {Ö}{{\"O}}1
           {ä}{{\"a}}1
           {Ä}{{\"A}}1
           {ü}{{\"u}}1
           {Ü}{{\"U}}1
           {\\\%}{{\%}}1,
  basicstyle=\small\ttfamily,                   % small code
  breaklines=true,                              % break long lines
  commentstyle=\itshape\color{green!50!black},  % comments are green
  keywordstyle=[1]\color{blue!80!black},        % instructions are blue
  keywordstyle=[2]\color{orange!80!black},      % sections/other directives are orange
  keywordstyle=[3]\color{red!50!black},         % registers are red
  stringstyle=\color{mauve},                    % strings are from the telekom
  identifierstyle=\color{teal},                 % user declared addresses are teal
  frame=l,                                      % black line on the left side of code
  language=[RISC-V]Assembler,                   % all code is RISC-V
  tabsize=4,                                    % indent tabs with 4 spaces
  showstringspaces=false                        % do not replace spaces with weird underlines
}

\usepackage{comment}
\usepackage{enumitem}
%\usepackage{todonotes}
\usepackage[disable]{todonotes}
\presetkeys{todonotes}{inline}{}
\newcommand{\TODO}[1]{\todo[inline,color=orange]{TODO: #1}}
\usepackage[T1]{fontenc}
\geometry{margin=1.5cm}

\ConfigureHeadline{
    headline={title-name-id}
}

% Angaben fuer die Titelbox (Autor:innen, Semester, Kurztitel)
% ------------------------------------------------------------------
% Angaben fuer die Titelbox jedes Blatts.
% Kann auch direkt in ShareLaTeX angepasst werden.
% Der Titel selbst kommt aus der Testat-Vorlage (z. B. "Befehlskodierung").
% ------------------------------------------------------------------
\newcommand{\TestatKurztitel}{Testat RO SoSe27}   % erscheint in der Kopfzeile
\newcommand{\TestatAutoren}{Prof. Dr-Ing. Andreas Koch, David Volz, Florian Meisel}
\newcommand{\TestatSemester}{Sommersemester 2026}


% Zusaetzliche Pakete, die die Aufgaben dieses Testats brauchen
% (automatisch vom Generator erzeugt)
\input{zusatz}

% Nicht entfernen: \maketitle wird pro Matrikelnummer aufgerufen.
% Die Kopie ist noetig, weil KOMA-Script \maketitle nach dem ersten Aufruf abschaltet.
\let\BlattTitel\maketitle
